% !TEX root = scientific-computing.tex

% this needs to be at the top of each file with the \plot command and the word in the [] is the session for pythontex. 
\begin{jlcode}[prob]
using Plots, PGFPlotsX
pgfplotsx()
\end{jlcode}

\chapter{Using Random Numbers and \texorpdfstring{Probability\\ Models}{Probability Models}} \label{ch:prob-models}

To understand basic probability, often problems are examined that use combinatorics or counting techniques to solve them.  Consider 

\begin{quote}
A round table has 7 chairs around it.  Mary and her friend Alisha and 5 other people are given seats at the table in a random manner.  What is the probability that Mary and Alisha sit next to each other?
\end{quote}

After practicing with problems like this, the solutions can be found quite simply, but in this chapter we examine some ways to use simulation to find the solution.

\section{Rolling Dice}

A relatively-simple example is that of rolling dice.  As we saw in Chapter \ref{ch:random}, if we have a fair, six-sided die, the the probability of any of the numbers coming up is 1/6.  We can simulate that using random numbers in the following way:
\begin{jlblock}[prob]
S100 = rand(1:6,100)
\end{jlblock}
generated 100 numbers taken from the set $\{1,2,3,4,5,6\}$.  Running this command\footnote{don't forget that when running code with random numbers, you won't get the exact same results, but the spirit should be the same.} results in \jlc[prob]{print(S100[1:10])} for the first 10 elements. 

As the number of time we roll a die increases, the probability of any one number appearing gets closer to 1/6.  For example, to test this, try
\begin{jlblock}[prob]
p = count(a->a==3,rand(1:6,1_000))
\end{jlblock}
and then evaluating \verb!p/1_000! results in \jlc[prob]{print(p/1_000)} \printpythontex[verb], which is reasonably close to 1/6. 


\subsection{Exercise}

\begin{itemize}
\item Trying increasing the number of random numbers used in the above example.  You should notice that as that number gets larger, the resulting probability gets closer to 1/6. 
\item What does 
\begin{jlblock}[prob]
p = count(a -> a % 2 == 0, rand(1:6,1000))
\end{jlblock}
measure in the sense of a rolling a fair 6-sided die?  Does the result make sense?

\end{itemize}


\section{Flipping Coins}

Another probability model that arises often is that of flipping fair (and unfair) coins.  There are a number of ways to model this, however, we will use random boolean values.  For example, we can flip 100 coins by \jlb[prob]{rand(Bool,100)}.  (Note: the result looks like an array of 0s and 1s, but look at the type, \verb!Array{Bool,1}!, which says it is a 1-dimensional boolean array.)

We can determine the number of heads (when the result is 1), by
\begin{jlblock}[prob]
sum(rand(Bool,100))
\end{jlblock}
which returns \jlc[prob]{print(sum(rand(Bool,100)))} \printpythontex for me.  


\subsection{Flipping multiple coins}

Another simple example is to flip multiple coins and generally count the number of heads or tails seen.  Consider flipping 3-coins--say a penny, nickel and dime--and counting the number of heads.  We then do that 3-coin flip a larger number of times.  

We can do this with
\begin{jlblock}[prob]
coins3 = rand(Bool,100,3)
\end{jlblock}
and the top  10 lines of this is: \jlc[prob]{display(coins3[1:10,:])} \printpythontex[verbatim]

Each row contains the coin flips.  For example the first row has heads on all three.  The second row has a head and then two tails.  If we are interested in the sum of the number of heads, we can do this with the \verb!mapslices! functions seen in Chapter \ref{ch:arrays}:
\begin{jlblock}[prob]
num_heads = mapslices(sum,coins3;dims=[2])
\end{jlblock}
The first few elements of the array is \verb![3; 1; 2; 2; 3; 2; 1; 1 ...]! We can plot the results of this with the function
%
\begin{jlblock}[prob]
using Plots
histogram(num_heads,nbins=4)
\end{jlblock}
and the result is
\begin{center}
\pgfplotsset{scale=0.8}
\plot{plots/prob-models/num-heads.tex}{prob}
\end{center}


\section{Rolling 2 dice}

How do we handle the rolling of two dice? Here's an array with each row
having 2 dice.

\begin{jlblock}[prob]
dice2=rand(1:6,100,2)
\end{jlblock}
%
and then to find the sum of the dice:

\begin{jlblock}[prob]
dicesum = mapslices(sum,dice2;dims=[2])
\end{jlblock}
%
which sums along the rows. This is 100 rolls of 2 dice with the sum
recorded. First, to get an idea of the distribution of the dice sums,
let's plot the histogram
%
\begin{jlblock}[prob]
histogram(dicesum,nbins=11,xticks=2:2:12,legend=false)
\end{jlblock}
%
where \texttt{nbins} is the number of bins and in this case, since it
runs from 2 to 12, there's a total of 11 and this generates the plot: 
\begin{center}
\pgfplotsset{scale=0.7}
\plot{plots/prob-models/two-dice.tex}{prob}
\end{center}

We could find the number of 2's 3's, etc. using the \texttt{mapslices}
function as above, however there is a nice way to do this use the
\texttt{StatsBase} package. You may need to add the package and then
load it with

\begin{jlblock}[prob]
using StatsBase
\end{jlblock}

and then the \texttt{counts} function (
\href{https://juliastats.org/StatsBase.jl/latest/counts/#StatsBase.counts}{Read
the online documentation}) can be used:

\begin{jlblock}[prob]
counts(dicesum,2:12)
\end{jlblock}
%
returns a vector of how many of the dice sum fall into each number. When I ran this, I got

\jlc[prob]{print(counts(dicesum,2:12))}\printpythontex[verbatim] 

and a nicer way to plot the histogram is
\begin{jlblock}[prob]
bar(2:12,counts(dicesum,2:12),xticks=2:2:12,legend=false)
\end{jlblock}
\begin{center}
\pgfplotsset{scale=0.7}
\plot{plots/prob-models/bar.tex}{prob}
\end{center}

\subsection{Exercise}

\begin{enumerate}
\item Change the code above to use 10,000 dice rolls. Estimate the probability
that you

\begin{enumerate}
\item roll a 7.
\item roll a 10 or greater.
\item roll an even number.
\end{enumerate}

\item Also, plot a the histogram using 1000 dice rolls.
\end{enumerate}

\section{Other Probability Models}

Let's return to the problem that problem at the beginning of this section.  We will solve this problem using the random modeling in this chapter. 

First let's first store the names of the people at the take as an array
\begin{jlblock}[prob]
names = ["Alisha", "Mary", "person1", "person2", "person3", "person4", "person5"]
\end{jlblock}
where we use generic names for the other 5 people.  We will take random permutations of this array below and the determine if Alisha and Mary are next to each other. 

Before we find a random permutation, let's write a function that takes an array of names and returns \verb!true! if they are next to each other and \verb!false! if not. 

\begin{jlblock}[prob]
function nextToEachOther(names::Array{String,1})

# return true or false
end
\end{jlblock}

Like many aspects of coding there are many ways of doing this.  Here's what we will do:

\begin{enumerate}
\item Find the position in the array where Alisha is sitting.  This would be the place at the table she is in.
\item Find the position in the array where Mary is sitting. 
\item Determine if the two numbers are next to each other.  Don't forget that this could include positions 1 and 7. 
\end{enumerate}

The following is a relatively simple function is
\begin{jlblock}[prob]
function nextToEachOther(names::Array{String,1})
  a = findfirst(name -> name=="Alisha",names)
  m = findfirst(name -> name=="Mary",names)
  abs(a-m) == 1 || abs(a-m) == length(names)-1
end
\end{jlblock}
%
where the \verb!findfirst! function returns the index of the array where the function is true (that is where the two friends are sitting).  To test this:

\begin{jlblock}[prob]
nextToEachOther(["Alisha", "Mary", "person1", "person2", "person3", "person4", "person5"])
\end{jlblock} 
%
returns \jlc[prob]{print(nextToEachOther(["Alisha", "Mary", "person1", "person2", "person3", "person4", "person5"]))} \printpythontex

\begin{jlblock}[prob]
nextToEachOther(["Alisha", "person1", "Mary", "person2", "person3", "person4", "person5"])
\end{jlblock}
% 
 returns \jlc[prob]{print(nextToEachOther(["Alisha","person1","Mary","person2","person3","person4","person5"]))} \printpythontex

\begin{jlblock}[prob]
nextToEachOther(["Alisha","person1", "person2", "person3", "person4", "person5", "Mary"])
\end{jlblock}
%
 returns \jlc[prob]{print(nextToEachOther(["Alisha","person1","person2","person3","person4","person5", "Mary"]))} \printpythontex

\subsection{Random Permutations}

We now return to the problem to study the probability.  The \verb!shuffle! command in the \verb!Random! package\footnote{The Random package is a built-in package and doesn't need to be added, but just enter \jlb[prob]{using Random}} takes any array and shuffles (permutes) the contents in a random manner.  For example:
\begin{jlblock}[prob]
shuffle(names)
\end{jlblock}
returns \jlc[prob]{print(shuffle(names))} \printpythontex[verbatim][frame=single]

Then we can test if they are sitting next to each other.  The follow repeats this a large number of times:

\begin{jlblock}[prob]
function numTimes(trials::Integer)
  s = 0  # keeps track of how many times they sit next to each other
  for i=1:trials
    if nextToEachOther(shuffle(names))
      s += 1
    end
  end
  s/trials
end
\end{jlblock}
and then the fraction of times is \jlc[prob]{print(numTimes(10_000))} \printpythontex[verb]. 

The true value of the probability can be found in the following way. Mary sits in one of the seven chairs.  Alisha has a equal chance of sitting in one of the remaining 6 chairs.  Two of the chairs are next two Mary, so the probability is 2/6 or 1/3.  The result we see above is close to this value. 


